%------------------------------------------------------------------------------%
\begin{question}
\begin{columns}
\column{0.5\linewidth}
  Calcule o vetor tangente à curva
  \[
    x = t\cos(t)
  \]

  \[
    y = t\sin(t)
  \]

  \[
    0 \leq t < \infty
  \]
  quando \(t=\frac{\pi}{2}\)
\column{0.5\linewidth}
  \pause
  ~\\[5mm]
  \includegraphics{B_03_exemplo_3_curva}
\end{columns}
\end{question}

%------------------------------------------------------------------------------%
\begin{resolution}{Solução}
  \begin{columns}
    \column{0.5\linewidth}
    \onslide<+->{Derivada da coordenada $x$}
    \begin{align*}
      \onslide<+->{x'(t) & = \frac{dx}{dt}                \\}
      \onslide<+->{      & = \frac{d}{dt}t\cos(t)         \\}
      \onslide<+->{      & = 1\times\cos(t) + t(-\sin(t)) \\}
      \onslide<+->{      & = \cos(t) - t\sin(t)             }
    \end{align*}

    \column{0.5\linewidth}
    \onslide<+->{Derivada da coordenada $y$}
    \begin{align*}
      \onslide<+->{y'(t) & = \frac{dy}{dt}              \\}
      \onslide<+->{      & = \frac{d}{dt}t\sin(t)       \\}
      \onslide<+->{      & = 1\times \sin(t) + t\cos(t) \\}
      \onslide<+->{      & = \sin(t) + t\cos(t)           }
    \end{align*}
  \end{columns}
\end{resolution}

%------------------------------------------------------------------------------%
\begin{resolution}{Solução}
  \onslide<+->{Avaliando as funções derivada em \(t=\frac{\pi}{2}\)}
  \[
    \onslide<+->{x'\left(\frac{\pi}{2}\right)}
    \onslide<+->{   = \left(\cos(t) - t\sin(t)\right) \evalat{\frac{\pi}{2}}{}     }
    \onslide<+->{   = \cos\left(\frac{\pi}{2}\right) - \frac{\pi}{2}\sin\left(\frac{\pi}{2}\right)}
    \onslide<+->{   = 0 - \frac{\pi}{2} \times 1}
    \onslide<+->{   = -\frac{\pi}{2}}
  \]
  \[
    \onslide<+->{y'\left(\frac{\pi}{2}\right)}
    \onslide<+->{   = \left(\sin(t) + t\cos(t)\right) \evalat{\frac{\pi}{2}}{}     }
    \onslide<+->{   = \sin\left(\frac{\pi}{2}\right) + \frac{\pi}{2}\cos\left(\frac{\pi}{2}\right)}
    \onslide<+->{   = 1 + \frac{\pi}{2} \times 0}
    \onslide<+->{   = 1}
  \]

  \bigskip
  \onslide<+->{O vetor tangente é
    \(
    \begin{bmatrix}
      \sfrac{-\pi}{2} \\ 1
    \end{bmatrix}
    \)}
\end{resolution}

%------------------------------------------------------------------------------%
\begin{resolution}{Solução}
  \centering
  \includegraphics{B_03_exemplo_3_vetor}
\end{resolution}

%------------------------------------------------------------------------------%